\documentclass[11pt,a4paper]{article}

% --- THIẾT LẬP LỀ TRANG & KHỔ GIẤY ---
\usepackage[top=1.8cm,bottom=1.8cm,left=1.8cm,right=1.8cm]{geometry}

% --- HỖ TRỢ ĐA DẠNG TRÌNH BIÊN DỊCH (XeLaTeX / pdfLaTeX / LuaLaTeX) ---
\usepackage{iftex}
\ifXeTeX
  \usepackage{fontspec}
  \usepackage[vietnamese]{babel}
  \defaultfontfeatures{Ligatures=TeX}
\else
  \usepackage[utf8]{inputenc}
  \usepackage[T5,T1]{fontenc}
  \usepackage[vietnamese]{babel}
  \usepackage{mathptmx}
\fi

% --- CÁC GÓI TOÁN HỌC CHUẨN AMS ---
\usepackage{amsmath}
\usepackage{amssymb}
\usepackage{amsfonts}
\usepackage{mathtools}
\usepackage{bm}

% --- BẢNG BIỂU, DANH SÁCH & TRANG TRÍ ---
\usepackage{tabularx}
\usepackage{booktabs}
\usepackage{array}
\usepackage{enumitem}
\usepackage{multicol}
\usepackage{fancyhdr}
\usepackage{xcolor}

% --- THIẾT LẬP HEADER & FOOTER ---
\pagestyle{fancy}
\fancyhf{}
\renewcommand{\headrulewidth}{0pt}
\lfoot{\textit{\small Bộ đề kiểm tra Toán THCS KNTT - Mã: c1\_6\_bai4\_test1}}
\rfoot{\textit{\small Trang \thepage}}

\begin{document}

% ==============================================================================
% PHẦN ĐẦU ĐỀ THI (HEADER TIÊU CHUẨN BỘ GD&ĐT)
% ==============================================================================
\noindent
\begin{tabular*}{\textwidth}{@{\extracolsep{\fill}}lr@{}}
\textbf{SỞ GD\&ĐT ........................} & \textbf{ĐỀ KIỂM TRA ĐỊNH KỲ NĂM HỌC 2025 - 2026} \\
\textbf{TRƯỜNG THCS ........................} & \textbf{Môn: TOÁN - LỚP 6} \\
\textit{(Đề thi chuẩn ma trận 2025)} & \textit{Thời gian làm bài: 45 phút (không kể giao đề)}
\end{tabular*}

\vspace{6pt}
\begin{center}
{\Large\bfseries Toán 6 (KNTT) - Bài 4: Phép cộng và phép trừ số tự nhiên - Bài kiểm tra số 1}\\[4pt]
{\bfseries [ MỨC ĐỘ 1: CƠ BẢN - Mã đề: \texttt{c1\_6\_bai4\_test1} ]}
\end{center}

\vspace{4pt}
\noindent
\rule{\textwidth}{0.8pt}\\[3pt]
Họ và tên thí sinh: \dotfill\ Lớp: \dots\dots\ SBD: \dots\dots\\[-6pt]
\rule{\textwidth}{0.8pt}

\vspace{10pt}

% ==============================================================================
% PHẦN I: TRẮC NGHIỆM NHIỀU PHƯƠNG ÁN (3,0 ĐIỂM)
% ==============================================================================
\noindent
\textbf{\large PHẦN I. Câu trắc nghiệm nhiều phương án lựa chọn (3,0 điểm)}\\
\textit{\small Thí sinh trả lời từ câu 1 đến câu 12. Mỗi câu hỏi thí sinh chỉ chọn một phương án đúng nhất.}
\vspace{6pt}
\noindent\textbf{Câu 1.} [Nhận biết] Trong phép tính cộng$a + b = c$,$c$được gọi là gì?\par\vspace{2pt}
\noindent\begin{tabular*}{\textwidth}{@{\extracolsep{\fill}}llll@{}}
\textbf{A.} Tổng & \textbf{B.} Hiệu & \textbf{C.} Số hạng & \textbf{D.} Số bị trừ
\end{tabular*}\par\vspace{5pt}

\noindent\textbf{Câu 2.} [Nhận biết] Trong phép trừ$a - b = c$($a \ge b$),$a$được gọi là gì?\par\vspace{2pt}
\noindent\begin{tabular*}{\textwidth}{@{\extracolsep{\fill}}llll@{}}
\textbf{A.} Số trừ & \textbf{B.} Số bị trừ & \textbf{C.} Hiệu & \textbf{D.} Số hạng
\end{tabular*}\par\vspace{5pt}

\noindent\textbf{Câu 3.} [Nhận biết] Với mọi số tự nhiên$a$, kết quả của phép tính$a + 0$là:\par\vspace{2pt}
\noindent\begin{tabular*}{\textwidth}{@{\extracolsep{\fill}}llll@{}}
\textbf{A.} $0$ & \textbf{B.} $1$ & \textbf{C.} $a$ & \textbf{D.} $2a$
\end{tabular*}\par\vspace{5pt}

\noindent\textbf{Câu 4.} [Nhận biết] Tính chất giao hoán của phép cộng số tự nhiên được biểu diễn bởi công thức nào dưới đây?\par\vspace{2pt}
\noindent\begin{tabular*}{\textwidth}{@{\extracolsep{\fill}}p{0.48\textwidth}p{0.48\textwidth}@{}}
\textbf{A.} $a + b = b + a$ & \textbf{B.} $(a + b) + c = a + (b + c)$ \\
\textbf{C.} $a + 0 = a$ & \textbf{D.} $a - b = b - a$
\end{tabular*}\par\vspace{5pt}

\noindent\textbf{Câu 5.} [Nhận biết] Điều kiện để thực hiện được phép trừ$a - b$trong tập hợp số tự nhiên$\mathbb{N}$là gì?\par\vspace{2pt}
\noindent\begin{tabular*}{\textwidth}{@{\extracolsep{\fill}}llll@{}}
\textbf{A.} $a < b$ & \textbf{B.} $a \le b$ & \textbf{C.} $a > b$ & \textbf{D.} $a \ge b$
\end{tabular*}\par\vspace{5pt}

\noindent\textbf{Câu 6.} [Nhận biết] Kết quả của phép tính$47 + 0$bằng bao nhiêu?\par\vspace{2pt}
\noindent\begin{tabular*}{\textwidth}{@{\extracolsep{\fill}}llll@{}}
\textbf{A.} $0$ & \textbf{B.} $47$ & \textbf{C.} $470$ & \textbf{D.} $48$
\end{tabular*}\par\vspace{5pt}

\noindent\textbf{Câu 7.} [Nhận biết] Trong phép trừ$85 - 32 = 53$, số$32$được gọi là:\par\vspace{2pt}
\noindent\begin{tabular*}{\textwidth}{@{\extracolsep{\fill}}llll@{}}
\textbf{A.} Số bị trừ & \textbf{B.} Hiệu & \textbf{C.} Số trừ & \textbf{D.} Tổng
\end{tabular*}\par\vspace{5pt}

\noindent\textbf{Câu 8.} [Nhận biết] Công thức$(a + b) + c = a + (b + c)$thể hiện tính chất nào của phép cộng?\par\vspace{2pt}
\noindent\begin{tabular*}{\textwidth}{@{\extracolsep{\fill}}p{0.48\textwidth}p{0.48\textwidth}@{}}
\textbf{A.} Tính chất giao hoán & \textbf{B.} Tính chất kết hợp \\
\textbf{C.} Tính chất phân phối & \textbf{D.} Tính chất cộng với số$0$
\end{tabular*}\par\vspace{5pt}

\noindent\textbf{Câu 9.} [Thông hiểu] Kết quả của phép tính$135 + 246$là:\par\vspace{2pt}
\noindent\begin{tabular*}{\textwidth}{@{\extracolsep{\fill}}llll@{}}
\textbf{A.} $371$ & \textbf{B.} $381$ & \textbf{C.} $380$ & \textbf{D.} $281$
\end{tabular*}\par\vspace{5pt}

\noindent\textbf{Câu 10.} [Thông hiểu] Kết quả của phép tính$452 - 128$là:\par\vspace{2pt}
\noindent\begin{tabular*}{\textwidth}{@{\extracolsep{\fill}}llll@{}}
\textbf{A.} $324$ & \textbf{B.} $334$ & \textbf{C.} $326$ & \textbf{D.} $224$
\end{tabular*}\par\vspace{5pt}

\noindent\textbf{Câu 11.} [Thông hiểu] Tìm số tự nhiên$x$, biết$x - 15 = 35$.\par\vspace{2pt}
\noindent\begin{tabular*}{\textwidth}{@{\extracolsep{\fill}}llll@{}}
\textbf{A.} $x = 20$ & \textbf{B.} $x = 30$ & \textbf{C.} $x = 50$ & \textbf{D.} $x = 40$
\end{tabular*}\par\vspace{5pt}

\noindent\textbf{Câu 12.} [Thông hiểu] Tìm số tự nhiên$x$, biết$120 + x = 180$.\par\vspace{2pt}
\noindent\begin{tabular*}{\textwidth}{@{\extracolsep{\fill}}llll@{}}
\textbf{A.} $x = 60$ & \textbf{B.} $x = 300$ & \textbf{C.} $x = 50$ & \textbf{D.} $x = 70$
\end{tabular*}\par\vspace{5pt}


% ==============================================================================
% PHẦN II: TRẮC NGHIỆM ĐÚNG / SAI (2,0 ĐIỂM)
% ==============================================================================
\vspace{6pt}
\noindent
\textbf{\large PHẦN II. Câu trắc nghiệm đúng sai (2,0 điểm)}\\
\textit{\small Thí sinh trả lời từ câu 1 đến câu 2. Trong mỗi ý a), b), c), d) ở mỗi câu, thí sinh chọn đúng hoặc sai.}\\
\textit{\small (Quy tắc tính điểm: Đúng 1 ý được 0,1đ; đúng 2 ý được 0,25đ; đúng 3 ý được 0,5đ; đúng 4 ý được 1,0đ).}
\vspace{6pt}
\noindent\textbf{Câu 1.} Xét tính đúng hoặc sai của các khẳng định sau về phép cộng và phép trừ số tự nhiên:\par\vspace{2pt}
\begin{enumerate}[label=\textbf{\alph*)}, leftmargin=1.2cm, itemsep=2pt, topsep=2pt]
  \item \textnormal{[Thông hiểu]} Phép tính$125 + 75$cho kết quả bằng$200$.
  \item \textnormal{[Thông hiểu]} Trong tập hợp số tự nhiên$\mathbb{N}$, luôn thực hiện được phép trừ$35 - 50$.
  \item \textnormal{[Vận dụng thấp]} Áp dụng tính chất giao hoán và kết hợp, ta có:$37 + 54 + 63 = (37 + 63) + 54 = 154$.
  \item \textnormal{[Vận dụng thấp]} Nếu$x$là số tự nhiên thỏa mãn$x + 28 = 28$thì$x = 1$.
\end{enumerate}
\vspace{6pt}

\noindent\textbf{Câu 2.} Một cửa hàng gạo ngày thứ nhất bán được$150\text{ kg}$gạo. Ngày thứ hai bán được nhiều hơn ngày thứ nhất$30\text{ kg}$gạo.\par\vspace{2pt}
\begin{enumerate}[label=\textbf{\alph*)}, leftmargin=1.2cm, itemsep=2pt, topsep=2pt]
  \item \textnormal{[Thông hiểu]} Số gạo cửa hàng bán được trong ngày thứ hai là$180\text{ kg}$.
  \item \textnormal{[Thông hiểu]} Tổng số gạo bán được trong cả hai ngày là$300\text{ kg}$.
  \item \textnormal{[Vận dụng thấp]} Nếu ban đầu cửa hàng có$400\text{ kg}$gạo thì sau hai ngày bán, cửa hàng còn lại$70\text{ kg}$gạo.
  \item \textnormal{[Vận dụng thấp]} Để bán hết số gạo còn lại trong ngày thứ ba thì ngày thứ ba cửa hàng phải bán$80\text{ kg}$gạo.
\end{enumerate}
\vspace{6pt}


% ==============================================================================
% PHẦN III: TRẮC NGHIỆM TRẢ LỜI NGẮN (2,0 ĐIỂM)
% ==============================================================================
\vspace{6pt}
\noindent
\textbf{\large PHẦN III. Câu trắc nghiệm trả lời ngắn (2,0 điểm)}\\
\textit{\small Thí sinh trả lời từ câu 1 đến câu 4. Mỗi câu trả lời đúng thí sinh được 0,5 điểm.}
\vspace{6pt}
\noindent\textbf{Câu 1.} [Thông hiểu] Tính giá trị của biểu thức:$148 + 52$.\par\vspace{2pt}
\noindent\textit{Đáp số:} \dotfill\par\vspace{6pt}

\noindent\textbf{Câu 2.} [Thông hiểu] Tính giá trị của biểu thức:$365 - 165$.\par\vspace{2pt}
\noindent\textit{Đáp số:} \dotfill\par\vspace{6pt}

\noindent\textbf{Câu 3.} [Thông hiểu] Tìm số tự nhiên$x$thỏa mãn:$x + 45 = 100$.\par\vspace{2pt}
\noindent\textit{Đáp số:} \dotfill\par\vspace{6pt}

\noindent\textbf{Câu 4.} [Thông hiểu] Một lớp học có$42$học sinh, trong đó có$23$học sinh nữ. Hỏi lớp học đó có bao nhiêu học sinh nam?\par\vspace{2pt}
\noindent\textit{Đáp số:} \dotfill\par\vspace{6pt}


% ==============================================================================
% PHẦN IV: TỰ LUẬN (3,0 ĐIỂM)
% ==============================================================================
\vspace{6pt}
\noindent
\textbf{\large PHẦN IV. Tự luận (3,0 điểm)}\\
\textit{\small Thí sinh trình bày chi tiết các bước giải từ câu 1 đến câu 3 vào giấy làm bài.}
\vspace{6pt}
\noindent\textbf{Bài 1 (1,0 điểm).} [Thông hiểu] Tính bằng cách hợp lý nhất:\par\noindent a)$48 + 125 + 52$b)$168 + 245 + 32 + 55$\par\vspace{8pt}

\noindent\textbf{Bài 2 (1,0 điểm).} [Thông hiểu] Tìm số tự nhiên$x$, biết:\par\noindent a)$x - 36 = 64$b)$215 - x = 105$\par\vspace{8pt}

\noindent\textbf{Bài 3 (1,0 điểm).} [Thông hiểu] Mẹ đưa cho Nam$100\,000$đồng đi mua đồ dùng học tập. Nam mua một cuốn sách tham khảo hết$45\,000$đồng và một bộ thước kẻ hết$15\,000$đồng. Hỏi Nam còn lại bao nhiêu tiền mang về cho mẹ?\par\vspace{8pt}


\begin{center}
\textbf{---------- HẾT ----------}\\
\textit{\small Cán bộ coi thi không giải thích gì thêm.}
\end{center}

\newpage
% ==============================================================================
% HƯỚNG DẪN CHẤM VÀ ĐÁP ÁN CHI TIẾT
% ==============================================================================
\begin{center}
{\Large\bfseries HƯỚNG DẪN CHẤM VÀ ĐÁP ÁN CHI TIẾT}\\[4pt]
{\large\bfseries Toán 6 (KNTT) - Bài 4: Phép cộng và phép trừ số tự nhiên - Bài kiểm tra số 1}\\[2pt]
{\bfseries Mã đề: \texttt{c1\_6\_bai4\_test1}}
\end{center}

\vspace{10pt}
\noindent\textbf{\large 1. ĐÁP ÁN PHẦN I: TRẮC NGHIỆM NHIỀU PHƯƠNG ÁN (3,0 điểm)}
\vspace{4pt}
\begin{center}
\begin{tabular}{|c|c|c|c|c|c|c|c|c|c|c|c|}
\hline
\textbf{1} & \textbf{2} & \textbf{3} & \textbf{4} & \textbf{5} & \textbf{6} & \textbf{7} & \textbf{8} & \textbf{9} & \textbf{10} & \textbf{11} & \textbf{12} \\
\hline
\textbf{A} & \textbf{B} & \textbf{C} & \textbf{A} & \textbf{D} & \textbf{B} & \textbf{C} & \textbf{B} & \textbf{B} & \textbf{A} & \textbf{C} & \textbf{A} \\
\hline
\end{tabular}
\end{center}
\vspace{6pt}


\noindent\textbf{\large 2. ĐÁP ÁN PHẦN II: TRẮC NGHIỆM ĐÚNG / SAI (2,0 điểm)}
\vspace{4pt}
\begin{center}
\begin{tabular}{|c|c|c|c|c|}
\hline
\textbf{Câu hỏi} & \textbf{Ý a)} & \textbf{Ý b)} & \textbf{Ý c)} & \textbf{Ý d)} \\
\hline
\textbf{Câu 1} & \textbf{Đúng} & \textbf{Sai} & \textbf{Đúng} & \textbf{Sai} \\
\hline
\textbf{Câu 2} & \textbf{Đúng} & \textbf{Sai} & \textbf{Đúng} & \textbf{Sai} \\
\hline
\end{tabular}
\end{center}
\vspace{6pt}


\noindent\textbf{\large 3. ĐÁP ÁN PHẦN III: TRẢ LỜI NGẮN (2,0 điểm)}
\vspace{4pt}
\begin{center}
\begin{tabular}{|c|m{11cm}|c|}
\hline
\textbf{Câu} & \textbf{Tóm tắt yêu cầu} & \textbf{Đáp số chính xác} \\
\hline
\textbf{Câu 1} & Tính giá trị của biểu thức:$148 + 52$. & \textbf{200} \\
\hline
\textbf{Câu 2} & Tính giá trị của biểu thức:$365 - 165$. & \textbf{200} \\
\hline
\textbf{Câu 3} & Tìm số tự nhiên$x$thỏa mãn:$x + 45 = 100$. & \textbf{55} \\
\hline
\textbf{Câu 4} & Một lớp học có$42$học sinh, trong đó có$23$học sinh nữ. Hỏi lớp học đó có ba... & \textbf{19} \\
\hline
\end{tabular}
\end{center}
\vspace{6pt}


\noindent\textbf{\large 4. LỜI GIẢI CHI TIẾT VÀ BIỂU ĐIỂM TỰ LUẬN (3,0 điểm)}
\vspace{4pt}
\noindent\textbf{\underline{Bài 1 (1,0 điểm):}}\par\vspace{3pt}
a) (0.5đ): Viết đúng nhóm số$(48 + 52) + 125$(0.25đ). Tính ra kết quả$100 + 125 = 225$(0.25đ).\par\noindent b) (0.5đ): Nhóm hợp lý$(168 + 32) + (245 + 55)$(0.25đ). Tính ra kết quả$200 + 300 = 500$(0.25đ).\par\vspace{8pt}

\noindent\textbf{\underline{Bài 2 (1,0 điểm):}}\par\vspace{3pt}
a) (0.5đ): Nêu được$x = 64 + 36$(0.25đ). Kết luận$x = 100$(0.25đ).\par\noindent b) (0.5đ): Nêu được$x = 215 - 105$(0.25đ). Kết luận$x = 110$(0.25đ).\par\vspace{8pt}

\noindent\textbf{\underline{Bài 3 (1,0 điểm):}}\par\vspace{3pt}
Bước 1 (0.5đ): Tính tổng số tiền Nam đã mua:$45\,000 + 15\,000 = 60\,000$(đồng).\par\noindent Bước 2 (0.5đ): Tính số tiền Nam còn lại mang về:$100\,000 - 60\,000 = 40\,000$(đồng) và ghi đúng đáp số.\par\vspace{8pt}

\end{document}
